\section{How proofs proceed}
\label{sec:proofs}

The dissertation gave few examples. This section gives three, in increasing
size. All are derivations in the calculus of Section~\ref{sec:calculus},
checked by Rocq.

\subsection{A store, with and without a fence}

For any invariant and either mode,
\[
  \mspec{J}{m}{x \fcell u}{\cseq{\cstore{x}{v}_i}{\cfence_i}}{x \fcell v}
\]
is derivable (\rocq{owned_derive_store_fence_cell}): \rn{store-any} gives
$x \fcell u \fseq x \fwrite{i} v$, \rn{frame-tm} carries that history
across $\cfence_i$, and a drained history collapses to its last write. This
is the store rule of ordinary separation logic. The same specification
without the fence is \emph{invalid} (\rocq{owned_plain_store_cell_invalid}):
the store may terminate with its write still buffered, and a cell assertion
has no buffered writes. The logic is the difference between these two.

\subsection{Reading a buffered write}

After $\cstore{x}{v}_i$ the thread holds $R \fseq x \fwrite{i} v$, where
$R$ describes the older state of $x$. The load axiom wants the loaded write
to be spatially separate from the rest, and an old cell cannot be separated
spatially from a newer write to the same address. The temporal frame rule
is what resolves this. The newest write alone satisfies the axiom's
precondition, with $\femp$ for $P$; the older history is a temporal frame:
\[
\inferrule*[right=frame-tm]
  {\mspec{J}{\modeov}{x \fwrite{i} v}{\cload{r}{x}_i}
         {x \fwrite{i} v \conj r = v}}
  {\mspec{J}{\modeov}{R \fseq x \fwrite{i} v}{\cload{r}{x}_i}
         {R \fseq (x \fwrite{i} v \conj r = v)}}
\]
(\rocq{owned_derive_load_after_history}). Store-to-load forwarding, the
behaviour most characteristic of the machine, is thus an instance of the
rule that most distinguishes the logic, and the reason it was worth
keeping.

\subsection{A shared container}

One cell $x$ is shared through a flag $f$. With
$X = \exists v \st x \fcell v$, the invariant is
\[
  J \;=\; f \fcell 0 \;\disj\; (f \fcell 1 \fsep X).
\]
A thread on processor $1$ takes the item, writes to it privately, and
returns it:
\begin{align*}
 & \clock_1;\ \cload{r}{f}_1;\ \cassume{r = 1}_1;\ \cstore{f}{0}_1;\
   \cunlock_1; \\
 & \cstore{x}{9}_1; \\
 & \clock_1;\ \cload{r}{f}_1;\ \cassume{r = 0}_1;\ \cstore{f}{1}_1;\
   \cunlock_1 .
\end{align*}
Meanwhile processor $2$ has an unrelated write $R = y \fwrite{2} 17$
buffered. The derived specification is
$\mspec{J}{\modedj}{R}{c}{R}$ (\rocq{owned_container_derived}). Writing
$H(a,v)$ for $(\exists u \st a \fcell u) \fseq a \fwrite{1} v$, the proof
passes through these assertions, each with $R$ as a spatial frame:
\begin{center}\small
\begin{tabular}{@{}ll@{}}
\toprule
After & The thread holds \\
\midrule
first load and assume & $\flock{1} \fsep (f \fcell 1 \fsep X)$ \\
$\cstore{f}{0}$ & $\flock{1} \fsep (H(f,0) \fsep X)$ \\
first unlock & $X$ \\
$\cstore{x}{9}$ & $H(x,9)$ \\
second lock & $H(x,9) \fseq ((\flock{1} \fsep J) \fseq \fbar{1})$ \\
second load and assume & $\flock{1} \fsep (f \fcell 0 \fsep x \fcell 9)$ \\
$\cstore{f}{1}$ & $\flock{1} \fsep (H(f,1) \fsep x \fcell 9)$ \\
second unlock & $\femp$ \\
\bottomrule
\end{tabular}
\end{center}
Three things are visible in it.

\emph{The foreign write costs one rule application.} $R$ concerns an
address the thread never touches, so a single \rn{frame-sp} around the
entire program carries it, across both releases.

\emph{The temporal frame rule carries the private write across the second
acquisition.} $H(x,9)$ is overlapping history, possibly still buffered.
\rn{frame-tm} places it before the acquired state, and the barrier that
\rn{lock} delivers then drains it, leaving $x \fcell 9$.

\emph{Releases use entailment, not framing.} The two \rn{unlock} premises
are $(H(f,0) \fsep X) \fseq \fbar{1} \sentails J \fsep X$ and
$(H(f,1) \fsep x \fcell 9) \fseq \fbar{1} \sentails J \fsep \femp$. The
first keeps the item and the second donates it.

The example is not vacuous. From a memory with the lock free, $f = 1$,
$x = 7$ and processor $2$'s write pending, there is a complete 24-step
execution in which that write is pending in every state
(\rocq{owned_container_foreign_pending_throughout}), and by adequacy every
execution from that memory is safe and, if it completes, restores the
invariant.

A second version replaces each $\cassume{}$ by a loop that releases the
lock and retries when the flag is wrong (\rocq{owned_retrying_container_derived}).
Its return loop has a branch for ``the container is already full'', and
that branch is derived from $\bexpf$: the thread holds $x$, so the pool
cannot.

\paragraph{The cost.}
The container's derivation applies \rn{cons} ten times and needs sixteen
entailments that are not identities, proved here semantically. Most are
rearrangements, a few are the laws about barriers and histories used
above, and one, at the second acquisition, takes fifty lines.
\pending{the same derivation with syntactic entailments; the Section 4.1
example; named locks and a producer/consumer client}
